\chapter*{Conclusion Générale}
\addcontentsline{toc}{chapter}{Conclusion Générale}
\markboth{Conclusion Générale}{Conclusion Générale}

%--------------------------------------------------------------
\section*{I. Bilan des contributions}
%--------------------------------------------------------------

\todo{%
  Reprendre les 7 modules et résumer ce qui a été livré dans chacun
  (2--3 lignes par module).
}

\begin{figure}[H]
  \centering
  \input{tikz/fig_conclusion_synthese}
  \caption{Synthèse — Architecture finale avec métriques obtenues}
  \label{fig:conclusion_synthese}
\end{figure}

\todo{%
  Reprendre la figure d'architecture globale (Fig.~\ref{fig:architecture_globale})
  et annoter chaque module avec ses métriques finales :
  WER STT, F1 NLU, latence E2E, nombre d'intents, taille KB, etc.
  C'est la figure que le jury photographiera pendant la soutenance.
}

%--------------------------------------------------------------
\section*{II. Difficultés rencontrées et solutions apportées}
%--------------------------------------------------------------

\begin{table}[H]
  \centering
  \caption{Difficultés rencontrées et solutions apportées}
  \label{tab:difficulties}
  \begin{tabularx}{\textwidth}{lXX}
    \toprule
    \textbf{Difficulté} & \textbf{Cause identifiée} & \textbf{Solution apportée} \\
    \midrule
    Binding RTP impossible sur IP Tailscale &
      IP \texttt{100.85.89.14} non bindable dans WSL &
      ESL\_HOST=127.0.0.1 \\
    Artéfacts audio &
      Resampling chunk par chunk &
      Accumulation int16 complète \\
    Double enregistrement CRM &
      Slot call\_logged non initialisé &
      Initialisation explicite à False \\
    Noms tunisiens mal prononcés &
      IPA espeak non fiable &
      Dictionnaire orthographique français (226 entrées) \\
    Race condition DuckDB &
      Initialisation ChromaDB asynchrone &
      Wrapper \_retry\_collection() avec back-off \\
    Port ESL 8021 occupé &
      svchost Windows &
      Migration vers port 8022 \\
    \bottomrule
  \end{tabularx}
\end{table}

%--------------------------------------------------------------
\section*{III. Limites actuelles}
%--------------------------------------------------------------

\todo{%
  Barge-in non implémenté, WER Darija dégradé, module émotionnel 3 classes,
  scalabilité mono-serveur.
}

%--------------------------------------------------------------
\section*{IV. Perspectives et travaux futurs}
%--------------------------------------------------------------

\todo{%
  1. Barge-in (couche RTP stabilisée)
  2. Fine-tuning STT corpus tunisien (LINAGORA arXiv 2504.02604)
  3. TTS dialectal tunisien
  4. Migration Celery complète pour le pipeline training
  5. Déploiement VPS Ubuntu
  6. Roadmap DevOps : k3s, GitLab CI, Prometheus/Grafana
}

